\section{结论}
本文围绕虚拟现实课程中的视觉分析与三维场景展示任务，完成了浏览器端人脸分析实验、dlib 人脸描述子训练示例分析，以及 Unity 森林 VR 场景构建三个部分的工作。报告首先分析了基于 \cmd{face-api.js} 的浏览器端人脸分析项目，说明了静态页面如何加载预训练模型，并完成表情识别、年龄性别估计和摄像头实时人脸追踪。随后，报告进一步分析 dlib 官方度量学习示例，从训练数据组织、mini-batch 采样、ResNet 描述子网络、\cmd{loss_metric} 损失函数和训练后阈值验证等方面解释了人脸描述子为什么可以通过距离进行比较。

在虚拟现实场景构建方面，本文使用 Unity VR 模板创建项目，并以森林为主题导入 Fantasy Forest Environment Free Sample 资源包。实验过程中遇到旧 Built-in 渲染管线资源导入 Unity 6 + URP 后材质变粉的问题。通过分析 Shader 与渲染管线的兼容关系，本文将资源包自带的 \cmd{StandardNoCulling.shader} 改写为 URP 可用的 HLSL Shader，并调整相关材质引用，使森林地形、草地、树叶和树干能够正常显示。最后，通过复制完整 XR Origin 到森林场景，完成了 VR 视角接入和运行展示。

通过这些实验可以看出，视觉智能和虚拟现实应用都不是单一算法或单一界面能够完成的任务。浏览器端人脸分析需要把模型权重、前端页面、媒体输入和 Canvas 可视化连接起来；dlib 度量学习示例需要理解数据组织、网络结构、损失函数和阈值验证之间的关系；Unity VR 场景则需要同时处理资源导入、渲染管线、Shader、材质、相机和 XR 组件。三部分内容虽然侧重点不同，但共同体现了课程实验中“模型能力、工程实现和展示环境”之间的联系。

本文仍存在一些不足。人脸分析部分主要使用预训练模型，没有重新训练和系统评估模型性能；dlib 示例分析以代码理解和运行记录为主，没有使用大规模数据集完成完整训练；Unity VR 场景已经完成基本展示，但交互内容仍较少，尚未加入任务引导、物体拾取、路径漫游或与人脸识别结果联动的展示模块。后续可以在现有森林 VR 场景基础上加入更多交互元素，并尝试把前两章的人脸识别结果以虚拟展板、提示信息或场景触发事件的形式呈现在 VR 环境中，使视觉识别与虚拟现实展示形成更完整的应用闭环。
